\section{引言}
\label{sec:introduction}

对于一个冻结的 TFM，每次预测都同时依赖查询样本和带标签上下文。TFM 在多个任务上进行预训练，并通过上下文学习对查询样本进行预测，而无需更新数据集特定参数
\citep{muller2023tabpfn,hollmann2025tabpfn,qu2025tabicl,ma2025tabdpt}。
这类模型可用于医疗、金融、科学与工业领域的分类和回归任务。即使模型和查询样本保持不变，改变上下文样本、特征或标签仍可能改变预测结果 \citep{abdullah2026contextsampling}。因此，模型的稳健性不仅取决于查询数据，也取决于上下文的使用方式。

一个任务可以通过压缩、检索、样本选择、特征或响应变换以及分布对齐等方式构造多种上下文
\citep{feuer2024tunetables,thomas2024localpfn,xu2025mixturepfn,
treerath2026tabaicl,hollmann2025tabpfn,khalifa2023dense}。这些构造在不同变化下能够提供互补证据，但不存在一种构造可以在所有变化下始终可靠。我们将冻结预测器在原始、未经修改的上下文上得到的输出称为 \emph{Base}。直接平均相互冲突的替代预测，可能使原本准确的 Base 预测变差。因此，需要一种查询级机制来判断这些替代预测是否为更新提供了足够一致的证据，并在证据较弱时将输出锚定到 Base。

\textbf{\method{}} 将这一决策实现为无需额外训练主干模型的查询自适应上下文集成。该方法保留原始上下文作为 Base，并通过样本组成、特征与响应视图以及分布对齐构造互补上下文。同一个冻结预测器通过统一的输出接口评估每个上下文。置信度与跨上下文一致性共同决定替代预测的聚合结果，任务特定门控则控制更新幅度：支持证据越强，聚合结果的影响越大；支持证据较弱时，预测仍锚定到 Base。上下文构造和预测融合均不使用真实目标标签。\method{}-shared 在预定义变化上采用固定上下文库，且不接收变化类别或强度；若推理前已知变化类别，\method{}-specific 会加入相应候选上下文，但仍不接收变化强度。该接口无需改变模型架构，便可迁移到 TabPFN、TabICL、TabDPT 和 LimiX 等不同 TFM。

我们在 OpenML 任务、全合成数据与自然分布偏移上评估 \method{}，覆盖五种受控上下文扰动和七种分布偏移。在四个冻结的 TFM（LimiX-2M、TabICL-v2、TabPFN-v3 和 TabDPT v1.2）上，\method{}-shared 在所有非零强度的 OpenML 扰动与偏移指标上均优于 Base，同时保持接近的干净数据得分。在 180 个对齐的合成实验单元上，该方法的平均成对稳定性胜率排名第一，分类和回归分别达到 $67.0\%$ 和 $87.6\%$。在给定变化类别先验时，\method{}-specific 进一步提升性能；同一接口也适用于常规预测器。这些结果表明，可配置上下文能够在无需重新训练的情况下提高推理时稳健性。
\FloatBarrier
